\section{ROS 2 Integration and Interface Contract}\label{sec:ros}

The estimator cores of Section~\ref{sec:cores} carry no middleware
dependency; they are made deployable by thin ROS~2 node packages
(\texttt{prism\_loc} for the \texttt{laser2d} and \texttt{ndt3d} MCL
backends, \texttt{prism\_loc\_fusion\_ros} for the \texttt{fusion3d}
ESKF backend) whose job is to translate messages, TF, and parameters
into core calls and back. We treat this boundary as a first-class
interface rather than glue code, following the conventions that
\texttt{robot\_localization}~\cite{moore2014generalized} and
\texttt{nav2\_amcl} within Nav2~\cite{macenski2020marathon}
established for drop-in localization: a REP-105 frame contract, a
\texttt{map}$\to$\texttt{odom} TF output, and \texttt{/initialpose}
seeding. A robot already configured for AMCL can therefore switch
backends by changing a launch file, not its TF tree or its consumers.

\subsection{Frames and \texttt{map}$\to$\texttt{odom} composition}
All three backends honor the REP-105 frame chain
\texttt{map}$\to$\texttt{odom}$\to$\texttt{base\_link}, exposed as the
\texttt{global\_frame}, \texttt{odom\_frame}, and \texttt{base\_frame}
parameters. Each backend estimates the robot body in the map,
$T_{\mathrm{map,base}}$, but---so as not to fight the odometry source
that owns \texttt{odom}$\to$\texttt{base\_link}---it broadcasts only
the correction transform $T_{\mathrm{map,odom}}$. Writing $\oplus$ for
pose composition and $(\cdot)^{-1}$ for inversion, the node computes
\begin{equation}
  T_{\mathrm{map,odom}} \;=\;
  T_{\mathrm{map,base}} \oplus T_{\mathrm{odom,base}}^{-1},
  \label{eq:mapodom}
\end{equation}
where $T_{\mathrm{odom,base}}$ is read from TF at the measurement
stamp. Equation~\eqref{eq:mapodom} guarantees the round-trip identity
$T_{\mathrm{map,odom}} \oplus T_{\mathrm{odom,base}} =
T_{\mathrm{map,base}}$, so downstream consumers recover the intended
body pose; both the planar core
(\texttt{compose}/\texttt{inverse} on $SE(2)$) and the fusion node
(\texttt{map\_base}$\cdot$\texttt{odom\_base}$^{-1}$ on $SE(3)$
Isometries) implement it, and each is checked by a ROS-free unit test
asserting the round trip to $10^{-9}$. If \texttt{odom}$\to$\texttt{base\_link}
is momentarily unavailable, the fusion node degrades gracefully to
publishing \texttt{map}$\to$\texttt{base\_link} directly rather than
emitting a wrong transform. The broadcast stamp is future-dated by
\texttt{transform\_tolerance} (default $0.1$\,s) to bound
extrapolation on the consumer side.

\subsection{Topic contract and QoS}
Table~\ref{tab:contract} lists the per-backend contract. Inputs and
outputs are uniform across backends wherever the sensing allows: every
backend outputs the same \texttt{map}$\to$\texttt{odom} TF and a
\texttt{PoseWithCovarianceStamped} on \texttt{\textasciitilde/pose},
and every backend accepts \texttt{/initialpose}. QoS profiles match
each publisher's expected semantics rather than using defaults: sensor
streams (\texttt{LaserScan}, \texttt{PointCloud2}, \texttt{Imu},
\texttt{NavSatFix}) are subscribed with \texttt{SensorDataQoS}
(best-effort, volatile), while the occupancy-grid map is subscribed
\texttt{transient\_local} and \texttt{reliable} so that the single
latched grid published once by \texttt{map\_server} is delivered even
though the subscriber joins later. A mismatch here is a common and
silent failure mode for new ROS users; fixing the profiles at the
interface removes it.

\begin{table}[t]
\centering
\caption{Per-backend topic contract. All backends additionally
subscribe \texttt{/initialpose} and the TF
\texttt{odom}$\to$\texttt{base\_link}, and broadcast
\texttt{map}$\to$\texttt{odom}.}
\label{tab:contract}
\small
\begin{tabular}{@{}lll@{}}
\toprule
Backend & Inputs & Outputs \\
\midrule
\texttt{laser2d} & \texttt{/scan}, \texttt{/map} & \texttt{\textasciitilde/pose}, \\
 & (\texttt{OccupancyGrid}) & \texttt{\textasciitilde/particle\_cloud} \\[2pt]
\texttt{ndt3d} & \texttt{/points}, \texttt{.pcd} & \texttt{\textasciitilde/pose}, \\
 & (file) & \texttt{\textasciitilde/particle\_cloud} \\[2pt]
\texttt{fusion3d} & \texttt{/points}, \texttt{/imu}, & \texttt{\textasciitilde/pose}, \\
 & \texttt{/gnss} (\texttt{NavSatFix}), \texttt{.pcd} & \texttt{\textasciitilde/odometry} \\
\bottomrule
\end{tabular}
\end{table}

\subsection{Initialization paths}
The interface offers one manual and one automatic seeding path per
backend. Manually, all backends accept a
\texttt{PoseWithCovarianceStamped} on \texttt{/initialpose} (the RViz
``2D Pose Estimate'' tool): the MCL backends spread a Gaussian particle
cloud at the supplied mean and covariance, and the fusion backend loads
the pose as its initial nominal state. Automatically, the two backends
differ because their sensing differs. The \texttt{laser2d} backend can
self-recover with no prior pose by running the branch-and-bound search
of Section~\ref{sec:global} (\texttt{try\_global\_localization}, or the
on-demand \texttt{\textasciitilde/global\_localization} service), which
re-seeds the filter from a single scan. The \texttt{fusion3d} backend
auto-initializes once it has both an attitude estimate---roll and pitch
recovered from the IMU specific-force vector, heading from
\texttt{initial\_yaw}---and a position fix from the first accepted GNSS
message, converted from WGS-84 to a local ENU frame; \texttt{/initialpose}
overrides this when GNSS is unavailable.

\subsection{Parameter surface}
Every tunable is a declared ROS~2 parameter with an in-code default,
documented in a single source of truth (\texttt{PARAMS.md}) generated
from the \texttt{declare\_parameter} calls. The surface deliberately
exposes the core's control knobs rather than hiding them: the
KLD-sampling bounds (\texttt{kld\_err}, \texttt{kld\_z}, and the
histogram bin sizes), the likelihood-field and laser-range parameters
(\texttt{z\_hit}, \texttt{sigma\_hit}, \texttt{laser\_min\_range},
\texttt{laser\_max\_range}), the branch-and-bound search window and
depth (\texttt{bbs\_linear\_window}, \texttt{bbs\_angular\_step},
\texttt{bbs\_max\_depth}, \texttt{bbs\_min\_score\_fraction}), and the
ESKF process noise and GNSS gates (\texttt{sigma\_acc},
\texttt{gnss\_min\_status}, \texttt{gnss\_max\_pos\_cov}). Because the
cores are configured purely through plain structs, each parameter maps
to exactly one field, and the shipped YAML files override only the
subset a deployment typically changes.

\subsection{Runtime diagnostics}
A localization node that silently does nothing is the dominant
first-run frustration; this release invests in making every stalled
precondition observable and self-explaining. Four mechanisms cooperate.
(i)~\emph{Fail-fast on a missing mandatory map}: the \texttt{ndt3d} and
\texttt{fusion3d} nodes abort construction with an actionable
\texttt{std::runtime\_error} if \texttt{map\_pcd\_path} is empty or
fails to load, rather than idling forever. (ii)~\emph{Throttled
waiting-state warnings}: whenever a callback cannot yet act---no map,
no initial pose, no filter initialization---it emits a rate-limited
\texttt{WARN} that names both the unmet precondition and the remedy
(e.g.\ ``waiting for an initial pose---publish to \texttt{/initialpose}
\ldots or enable \texttt{try\_global\_localization}''). (iii)~A
\emph{$10$\,s startup watchdog} reports, for each mandatory
input that has produced no message, the topic name and its live
publisher count (via \texttt{count\_publishers}), disambiguating ``no
publisher'' from ``publisher present but QoS-incompatible''; it cancels
itself once all inputs are seen. (iv)~\emph{Quantified rejection
logging}: the GNSS gate logs each rejected fix with the measured value
against the threshold it violated---status below \texttt{gnss\_min\_status},
NaN coordinates, or position variance above \texttt{gnss\_max\_pos\_cov}
---and the NDT update likewise reports non-convergence and fitness
versus \texttt{ndt\_max\_fitness}. These diagnostics carry no algorithmic
weight; they exist so that a misconfiguration is diagnosable from the
console log alone, consistent with the operational transparency Nav2
found necessary for long-term autonomy~\cite{macenski2020marathon}.
